% On-policy distillation: principle and two runs. English version.
% Build: cd .report && tectonic -X compile opd-report-en.tex
% Design: warmreport.sty. Figure data: python .report/make_data.py (from the
% repository root).
\documentclass[11pt,a4paper]{article}
\usepackage{warmreport}
\renewcommand{\wrshorttitle}{On-policy distillation from first principles}
\hypersetup{pdftitle={On-Policy Distillation from First Principles}}

\newcommand{\pis}{\pi_\theta}
\newcommand{\pit}{\pi_T}
\newcommand{\KL}{\mathrm{KL}}
\newcommand{\E}{\mathbb{E}}

% Figure labels (curves.tex).
\newcommand{\LblKL}{Teacher KL (nats/token)}
\newcommand{\LblReward}{Reward (share of tasks solved)}
\newcommand{\LblLength}{Episode length (tokens)}
\newcommand{\LblEntropy}{Student entropy (nats/token)}
\newcommand{\LblStep}{Training step, stages 1--4 back to back}
\newcommand{\LblQwenThree}{Qwen3-8B $\to$ Qwen3-1.7B}
\newcommand{\LblQwenThreeFive}{Qwen3.5-9B $\to$ Qwen3.5-2B}

\begin{document}

\wrtitle{Technical report · hands-on-policy-distillation}
  {On-Policy Distillation from First Principles}
  {Two small Qwen models learn to use tools}
  {October 2026 · Qwen3-8B $\to$ Qwen3-1.7B and Qwen3.5-9B $\to$ Qwen3.5-2B
  on one A100-80GB}

\begin{wrabstract}{Abstract}
We distil Qwen3-8B into Qwen3-1.7B and Qwen3.5-9B into Qwen3.5-2B with the
same code and hyperparameters on one A100-80GB GPU, in four stages:
grade-school math, a calculator tool, three tools over multi-turn episodes,
and distillation plus a task reward. Starting from what the student should
imitate, and where, we derive on-policy distillation as a policy gradient
whose per-token reward is the teacher's log-probability of the student's
token minus the student's own. Teacher KL falls in every distillation stage,
and task success rises although no reward is used: from 0.47 to 1.00 and
from 0.65 to 0.90 on calculator tasks, from 0.66 to 0.75 and from 0.75 to
0.96 on three-tool tasks. Adding the task reward through GRPO changes nothing
measurable, because 98.5\% and 89\% of sample groups receive identical
rewards. The Qwen3.5 run takes 2.5 times the GPU time, because its episodes
are 2--5 times longer.
\end{wrabstract}

\begin{tcbraster}[raster columns=4, raster equal height,
    raster column skip=3mm]
\begin{stat}{wr-terracotta}{0.47 $\to$ 1.00}
calculator tasks solved by the Qwen3 student, with no reward
\end{stat}
\begin{stat}{wr-ochre}{0.75 $\to$ 0.96}
three-tool tasks solved by the Qwen3.5 student, with no reward
\end{stat}
\begin{stat}{wr-wine}{98.5\%}
of Qwen3's GRPO groups had identical rewards: no signal
\end{stat}
\begin{stat}{wr-sage}{2.5$\times$}
the GPU time for Qwen3.5, whose episodes are 2--5 times longer
\end{stat}
\end{tcbraster}

\section{Introduction}

A small model that calls tools is cheap to serve, but it has to learn when
to call which tool, how to write the call, and what to do with the result.
There are three ways to teach it: supervised fine-tuning on episodes that a
larger teacher wrote (off-policy distillation;
\citealp{kim2016seqkd}); reinforcement learning from a task
reward, for example with GRPO \citep{shao2024deepseekmath}; and on-policy
distillation, in which the student writes the episodes and the teacher
grades every token \citep{agarwal2024gkd,gu2024minillm,lu2025onpolicy}. The
small Qwen3 models were themselves post-trained with on-policy distillation
from larger ones \citep{yang2025qwen3}.

This report explains from first principles why the third works
(Section~\ref{sec:principle}), how we trained (Section~\ref{sec:setup}), and
what two runs showed (Section~\ref{sec:results}). Both runs use one pipeline
(TRL's GRPO trainer with a teacher term added to the advantage) and the same
configs; only the models differ.

\section{The core idea from first principles}
\label{sec:principle}

\paragraph{Setting.}
A language model is a policy over tokens. The state $s_t$ is everything in
the context before the student's $t$-th token: the prompt (system prompt,
tool descriptions, task), the student's earlier tokens, and the results of
the tools it has called. The action is the next token $y_t$. The student
$\pis$ is trained; the teacher $\pit$ is frozen. Tool results are written by
the environment, not chosen by the student, so they are part of the state
and never an action.

We want the student to choose tokens, in every state, the way the teacher
would. The natural discrepancy between two next-token distributions is a KL
divergence, which leaves two choices: in which states to measure it, and in
which direction.

\paragraph{Which states: the student's own.}
Fine-tuning on teacher-written episodes trains the student on the states
the teacher visits. At test time the student generates, and visits its own.
After one token the teacher would not have written, the student is in a
state that no training example covered, where it is more likely to err
again: errors compound. In imitation learning with per-step error $\epsilon$
on the training distribution, the cost over $T$ steps can grow as
$O(\epsilon T^2)$ \citep{ross2010efficient}; training on states drawn from
the learner's own policy, labelled by an expert, brings it down to
$O(\epsilon T)$ \citep[DAgger;][]{ross2011dagger}. On-policy distillation is
this recipe with the teacher as the expert: the student writes the episodes,
and the teacher says what it would have done in every state the student
reached.

\paragraph{Which direction: reverse KL.}
The objective is the reverse KL, summed over the states of the student's own
episodes $\tau$:
\begin{equation}
  \mathcal{L}(\theta) = \E_{x}\,\E_{\tau\sim\pis}\Big[\sum_{t}
  \KL\big(\pis(\cdot\mid s_t)\,\big\|\,\pit(\cdot\mid s_t)\big)\Big].
  \label{eq:objective}
\end{equation}
By the chain rule of KL, this expected sum is the KL between the two models'
distributions over whole episodes; tool results are the same function of the
history under both models and cancel. Reverse KL suits distillation into a
small model for three reasons. (i)~It is an expectation under the student,
so it can be estimated from the student's samples and the teacher's
log-probabilities of those tokens: one teacher forward pass, no teacher
generation. (ii)~It is mode seeking: it punishes the student hard for
putting probability where the teacher puts almost none, and mildly for
leaving out some of the teacher's options. A student that cannot represent
everything a larger teacher does commits to one good behaviour instead of
averaging several \citep{gu2024minillm,lu2025onpolicy}. (iii)~It reaches its
minimum, zero, only when the student behaves like the teacher in the states
it visits, so nothing but imitation lowers it. Unlike a learned reward
model, it cannot be gamed; \citet{lu2025onpolicy} call it unhackable.

\paragraph{The gradient is a policy gradient.}
Fix a state $s$ and write $p=\pis(\cdot\mid s)$ and $q=\pit(\cdot\mid s)$,
with $a$ ranging over the vocabulary. Since $p(a)\nabla_\theta\log
p(a)=\nabla_\theta p(a)$ and $\sum_a\nabla_\theta p(a)=\nabla_\theta 1=0$,
\begin{align}
  \nabla_\theta \KL(p\,\|\,q)
  &= \sum_a \nabla_\theta p(a)\big[\log p(a)-\log q(a)\big]
    + \sum_a \nabla_\theta p(a) \nonumber\\
  &= \E_{a\sim p}\Big[\big(\log p(a)-\log q(a)\big)\,\nabla_\theta\log p(a)\Big].
  \label{eq:gradient}
\end{align}
Define, for the token $y_t$ the student actually sampled,
\begin{equation}
  k_t = \log\pis(y_t\mid s_t) - \log\pit(y_t\mid s_t).
  \label{eq:k}
\end{equation}
By \eqref{eq:gradient}, $k_t\,\nabla_\theta\log\pis(y_t\mid s_t)$ is an
unbiased one-sample estimate of the gradient of the KL at $s_t$. That is the
whole method:

\begin{keybox}{Core identity}
Minimizing the reverse KL \eqref{eq:objective} on the student's own samples
is REINFORCE \citep{williams1992reinforce}: minimize the surrogate loss
\begin{equation}
  \ell(\theta) = -\frac{1}{N}\sum_{t} A_t\,\log\pis(y_t\mid s_t),
  \qquad A_t = -k_t \ \ \text{(held constant)},
  \label{eq:loss}
\end{equation}
over the $N$ student tokens of a batch. The reward of token $t$ is
\begin{equation}
  r_t = \log\pit(y_t\mid s_t) - \log\pis(y_t\mid s_t).
  \label{eq:reward}
\end{equation}
\end{keybox}

The reward is the teacher's log-probability of what the student wrote plus
the student's own surprisal (an entropy bonus). Tokens the teacher likes
more than the student did are pushed up, tokens it likes less are pushed
down. Averaged over tokens, $k_t$ is also an unbiased estimate of the
per-token reverse KL; it is the teacher-KL curve in Figure~\ref{fig:curves}.

\begin{algobox}[label=alg:step]{One training step}
\begin{algorithmic}[1]
\Require student $\pis$ (LoRA), frozen teacher $\pit$, environment,
  prompts $x_1,\dots,x_P$
\For{each prompt $x_p$}
  \State sample $G$ episodes with $\pis$ (vLLM), running every tool call and
    appending its result
  \State $R_{p,i}\gets$ task reward of episode $i$;\quad
    $\hat A_{p,i}\gets(R_{p,i}-\bar R_p)/\sigma_p$
    \Comment{0 when $G=1$}
\EndFor
\State teacher forward pass over every episode:
  $k_t\gets\log\pis(y_t\mid s_t)-\log\pit(y_t\mid s_t)$
\State $A_t\gets c_R\,\hat A_{p,i}-c_T\,k_t$ for every student token $t$ of
  episode $(p,i)$
\State one AdamW step on
  $\ell(\theta)=-\frac{1}{N}\sum_t A_t\log\pis(y_t\mid s_t)$
  \Comment{$N$ student tokens}
\end{algorithmic}
\smallskip
{\small\color{wr-taupe}Stages differ only in the environment, the group size
$G$, and $(c_T, c_R)$.}
\end{algobox}

\begin{notebox}{Two approximations.}
First, $\mathcal{L}$ depends on $\theta$ also through the states, because
every token shapes the states that follow; the exact gradient of
\eqref{eq:objective} adds the future terms $\sum_{t'>t}k_{t'}$ to $A_t$.
Like \citet{lu2025onpolicy}, we drop them (a discount factor of zero): each
token answers only for its own disagreement, which lowers variance, and the
next batch trains on whatever states the updated student visits. Second, one
sampled token stands in for the sum over the vocabulary in
\eqref{eq:gradient}. GKD \citep{agarwal2024gkd} computes that sum exactly,
which needs both models' full next-token distributions at every position
(248{,}320 numbers per position for Qwen3.5); the sampled estimate needs one
number from the teacher.
\end{notebox}

\paragraph{Why a dense signal beats a sparse one.}
Compare RL from a binary task reward $R\in\{0,1\}$. GRPO samples a group of
$G$ episodes per prompt and gives every token of episode $i$ the same
advantage $\hat A_i=(R_i-\bar R)/\sigma_R$, normalized within the group. Two
things follow. An episode's reward carries at most one bit, spread over
hundreds of tokens: it says that the episode failed, not where. And if all
$G$ episodes of a prompt get the same reward, every advantage in the group
is zero. If the student solves a prompt with probability $p$, independently
across samples, a group carries a signal with probability
\begin{equation}
  P_{\text{signal}}(p) = 1 - p^G - (1-p)^G ,
  \label{eq:signal}
\end{equation}
which for $G=4$ is 0.875 at $p=\tfrac12$, 0.17 at $p=0.955$, and vanishes as
$p$ approaches 0 or 1. The teacher term instead gives one number per token,
$O(T)$ bits per episode rather than $O(1)$ \citep{lu2025onpolicy}, and it is
nonzero wherever student and teacher disagree, whether or not the episode
succeeded. It is also cheap: the teacher scores an episode in one forward
pass, parallel over positions, and never decodes.

\paragraph{Combining both signals.}
In general, the advantage of token $t$ of episode $i$ is
\begin{equation}
  A_t = c_R\,\hat A_i - c_T\,k_t ,
  \label{eq:advantage}
\end{equation}
with $c_R=0$ for pure distillation. The teacher term can at best make the
student as good as the teacher; only the task reward can push past the
teacher or fix its mistakes, and only through the groups that
\eqref{eq:signal} leaves informative. Algorithm~\ref{alg:step} is one
training step of every stage.

\begin{notebox}{Preconditions.}
$k_t$ compares two models' probabilities of the same token id, so student
and teacher need the same tokenizer: Qwen3 (151{,}936 tokens) and Qwen3.5
(248{,}320) cannot be mixed, which is why each family is distilled within
itself. Prompts and tool results enter $s_t$ but carry no loss. And the
teacher has to judge states it would never reach on its own, such as those
after a student's mistake.
\end{notebox}

\section{Training setup}
\label{sec:setup}

\paragraph{Models and hardware.}
Table~\ref{tab:models} lists the two teacher--student pairs. Qwen3.5
\citep{qwen2026qwen35} interleaves Gated DeltaNet layers
\citep{yang2025gateddeltanet}, a form of linear attention, with gated
attention layers, three to one. Both runs use non-thinking mode, LoRA \citep{hu2022lora} of rank 32 and $\alpha=32$ on
every linear layer of the student's language model (the Qwen3.5 checkpoints
also carry a vision tower, which these text-only tasks leave unused), and
one Colab A100-SXM4-80GB. The trainer is TRL's \texttt{GRPOTrainer}
\citep{vonwerra2020trl} with vLLM generation \citep{kwon2023vllm}. Our
subclass computes $k_t$ after each rollout batch and replaces GRPO's
advantage with \eqref{eq:advantage}; sampling, tool execution, masking and
the loss are TRL's.

\begin{table}[t]
\centering\small
\caption{The two teacher--student pairs. Everything else is identical.}
\label{tab:models}
\begin{tabular}{@{}lll@{}}
\toprule
\headrow & Qwen3 run & Qwen3.5 run \\
\midrule
Teacher (frozen, bf16) & Qwen3-8B & Qwen3.5-9B (language model only) \\
Student & Qwen3-1.7B & Qwen3.5-2B \\
Student layers & 28, all attention & 18 Gated DeltaNet + 6 gated attention \\
Vocabulary & 151{,}936 & 248{,}320 \\
Tool-call syntax & JSON & XML \\
Trainable LoRA parameters & 34.9\,M & 33.6\,M \\
\bottomrule
\end{tabular}
\end{table}

\paragraph{Curriculum.}
Table~\ref{tab:stages} lists the four stages. Stage~1 samples GSM8K questions
\citep{cobbe2021gsm8k}; their reference answers never enter the loss.
Stage~2's word problems have numbers large enough (a five-digit
times a four-digit number, say) that arithmetic in prose often goes wrong.
Stages~3 and~4 offer three tools on every task: a calculator, read-only SQL
over a randomly generated company database, and JMESPath queries over a
randomly generated store document. Their tasks need one tool (45\%), two
tools combined (40\%), or none (15\%, to be answered directly). The tasks of
stages 2--4 are generated, with different random seeds for training and
evaluation. Each stage keeps training the previous stage's LoRA adapter.

\begin{table}[t]
\centering\small
\caption{The curriculum, identical in both runs. Tokens: limit on the whole
episode after the prompt, tool results included. Batch: episodes per step,
prompts $\times$ group size $G$.}
\label{tab:stages}
\begin{tabular}{@{}lllcrrrcc@{}}
\toprule
\headrow Stage & Tasks & Tools & Turns & Tokens & Steps & Batch & LR
  & $(c_T, c_R)$ \\
\midrule
1 & GSM8K & none & 1 & 1{,}024 & 200 & $32\times1$ & $10^{-4}$ & (1, 0) \\
2 & big-number arithmetic & calculator & 4 & 2{,}048 & 150 & $32\times1$
  & $10^{-4}$ & (1, 0) \\
3 & one, two or no tools & all three & 6 & 3{,}072 & 200 & $32\times1$
  & $10^{-4}$ & (1, 0) \\
4 & as stage 3 & all three & 6 & 3{,}072 & 150 & $16\times4$
  & $5\cdot10^{-5}$ & (1, 1) \\
\bottomrule
\end{tabular}
\end{table}

\paragraph{Optimization and scoring.}
\begin{itemize}[leftmargin=*,itemsep=1pt,topsep=2pt]
\item Every optimizer step uses a fresh batch sampled by the current
  weights, so the PPO ratio is one, clipping ($\epsilon=0.2$ below, 0.28
  above) never binds, and the update is the gradient of \eqref{eq:loss},
  averaged over all student tokens of the batch (TRL's DAPO loss;
  \citealp{yu2025dapo}).
\item AdamW with a constant learning rate after 5 warm-up steps, gradient
  norm clipped at 1, no weight decay; micro-batches of 8 episodes with
  gradient checkpointing. A rate of $10^{-4}$ is about ten times a full
  fine-tuning rate, as recommended for LoRA \citep{schulman2025lora}.
\item Sampling: vLLM on the training GPU (30\% of its memory), temperature
  1, top-$p$ 1.
\item Scoring: one bf16 teacher forward pass per episode, through a fused
  output layer that returns the log-probabilities of the sampled tokens
  without materializing the full logits. The student's log-probabilities in
  $k_t$ are recomputed by the trainer, not taken from vLLM.
\item vLLM's kernels give slightly different log-probabilities from the
  trainer's. TRL weights each episode by
  $\exp\sum_t(\log\pis-\log\pi_{\text{vLLM}})$ and drops episodes whose
  weight exceeds 3; the mean weight per stage was 0.95--1.00.
\item Reward (stage-4 training, all evaluations): 1 if the last
  \verb|\boxed{}| of the final answer equals the reference (numerically, to
  $10^{-6}$), else 0; also 0 if the episode ran out of tokens or was still
  calling tools at the turn limit.
\item Evaluation every 25 steps on 256 held-out tasks (GSM8K test questions
  in stage~1), one sample per task at temperature 1, since TRL has no greedy
  evaluation.
\end{itemize}

\section{Results}
\label{sec:results}

Table~\ref{tab:results} gives the first and last evaluation of every stage;
Figure~\ref{fig:curves} shows all of them.

\begin{table}[t]
\centering\small
\caption{Evaluation at the start and end of each stage (256 held-out tasks,
temperature 1), and wall-clock time including evaluation.}
\label{tab:results}
\setlength{\tabcolsep}{4.5pt}
\begin{tabular}{@{}lcccccc@{}}
\toprule
\headrow & \multicolumn{3}{c}{Qwen3-8B $\to$ Qwen3-1.7B}
 & \multicolumn{3}{c}{Qwen3.5-9B $\to$ Qwen3.5-2B} \\
\headrow Stage & Teacher KL & Reward & Time & Teacher KL & Reward & Time \\
\midrule
1. GSM8K & 0.229 $\to$ 0.136 & 0.809 $\to$ 0.836 & 45 min
  & 0.218 $\to$ 0.168 & 0.656 $\to$ 0.781 & 85 min \\
2. Calculator & 0.328 $\to$ 0.057 & 0.473 $\to$ 1.000 & 18 min
  & 0.108 $\to$ 0.047 & 0.645 $\to$ 0.902 & 103 min \\
3. Three tools & 0.378 $\to$ 0.068 & 0.660 $\to$ 0.746 & 52 min
  & 0.163 $\to$ 0.034 & 0.754 $\to$ 0.957 & 109 min \\
4. + Task reward & 0.065 $\to$ 0.067 & 0.754 $\to$ 0.758 & 55 min
  & 0.034 $\to$ 0.025 & 0.945 $\to$ 0.953 & 130 min \\
\midrule
Total & & & 2.8 h & & & 7.1 h \\
\bottomrule
\end{tabular}
\end{table}

\begin{figure}[t]
\centering
% Evaluation curves of both runs, from data/*.dat (written by make_data.py).
% Shared by both reports: each defines the label macros \Lbl... before \input.
\pgfplotsset{
  run3/.style={color=wr-wine, mark=*, mark size=1.1pt, thick},
  run35/.style={color=wr-orange, mark=square*, mark size=1.1pt, thick},
}
% Stage numbers centred over each stage (x = 100, 275, 450, 625 of 700).
\newcommand{\stagenumbers}{%
  \node[font=\scriptsize, text=wr-taupe, anchor=north] at (rel axis cs:0.1429,1) {1};
  \node[font=\scriptsize, text=wr-taupe, anchor=north] at (rel axis cs:0.3929,1) {2};
  \node[font=\scriptsize, text=wr-taupe, anchor=north] at (rel axis cs:0.6429,1) {3};
  \node[font=\scriptsize, text=wr-taupe, anchor=north] at (rel axis cs:0.8929,1) {4};}
\begin{tikzpicture}
\begin{groupplot}[
  group style={group size=2 by 2, horizontal sep=1.5cm, vertical sep=1.45cm},
  width=0.5\textwidth, height=4.6cm,
  xmin=0, xmax=700, xtick={0,200,350,550,700},
  xmajorgrids,
  unbounded coords=jump,
  tick label style={font=\scriptsize},
  title style={font=\small\sffamily, text=wr-brown, yshift=-1ex},
  label style={font=\small},
]
\nextgroupplot[title={(a) \LblKL}, ymin=0, ymax=0.44,
  legend to name=curveslegend, legend columns=2,
  legend style={font=\small, draw=none,
    /tikz/every even column/.append style={column sep=1.5em}}]
\addplot[run3] table[x=step, y=kl] {data/qwen3.dat};
\addplot[run35] table[x=step, y=kl] {data/qwen35.dat};
\legend{\LblQwenThree, \LblQwenThreeFive}
\stagenumbers

\nextgroupplot[title={(b) \LblReward}, ymin=0.4, ymax=1.12]
\addplot[run3] table[x=step, y=reward] {data/qwen3.dat};
\addplot[run35] table[x=step, y=reward] {data/qwen35.dat};
\stagenumbers

\nextgroupplot[title={(c) \LblLength}, ymode=log, ymin=40, ymax=2000,
  xlabel={\LblStep}]
\addplot[run3] table[x=step, y=length] {data/qwen3.dat};
\addplot[run35] table[x=step, y=length] {data/qwen35.dat};

\nextgroupplot[title={(d) \LblEntropy}, ymode=log, ymin=5e-4, ymax=1,
  xlabel={\LblStep}]
\addplot[run3] table[x=step, y=entropy] {data/qwen3.dat};
\addplot[run35] table[x=step, y=entropy] {data/qwen35.dat};
\end{groupplot}
\end{tikzpicture}

\smallskip
\ref*{curveslegend}

\caption{Evaluations during the four stages, run back to back. Dashed lines
separate the stages; each stage opens with an evaluation of the previous
stage's weights on the new tasks. (a)~Mean of $k_t$ over student tokens.
(b)~Share of the 256 held-out tasks solved. (c)~Mean episode length after
the prompt (log scale). (d)~Mean entropy of the student's next-token
distribution (log scale).}
\label{fig:curves}
\end{figure}

\paragraph{Distillation does what the derivation says.}
Teacher KL falls in every distillation stage, by 23--83\%, and mostly within
the first 25--50 steps (Qwen3, stage~2: from 0.328 to 0.049 by step 25).

\paragraph{Task success follows without a reward.}
Stages 1--3 never see the reward, yet the evaluation reward rises in each:
by 0.03 and 0.13 on GSM8K, 0.53 and 0.26 on calculator tasks, 0.09 and 0.20
on three-tool tasks (Qwen3 and Qwen3.5). The students take up their
teachers' tool habits. In the first batch of stage~2, the Qwen3.5 student
called the calculator 4 times in 32 episodes and did the rest of the
arithmetic in prose; over the last 50 steps it called it 1.55 times per
episode (Qwen3: 15 times in 32 episodes, then 1.16 per episode). In
stage~3, the share of executed tool calls that returned an error fell from
21\% to 4\% (first against last five steps; Qwen3: from 14\% to 7\%).

\paragraph{Length, not accuracy, holds Qwen3.5 back.}
On GSM8K, 10--18\% of the Qwen3.5 student's answers run past the
1{,}024-token limit and score zero. They are not loops: after an awkward
intermediate result, such as a fractional count, the model second-guesses
itself. Of the answers that finish, 91\% are right at the end of stage~1
(Qwen3: 84\%). In stage~2, about 70\% of its failed training episodes in the
last 50 steps used all four turns re-checking a result with more calculator
calls.

\begin{findbox}{The task reward adds nothing measurable.}
In stage~4 the evaluation reward moves from 0.754 to 0.758 and from 0.945 to
0.953. A 256-task estimate near 0.95 has a standard error of 0.014, and the
same Qwen3.5 weights scored 0.957 at the end of stage~3.
Equation~\eqref{eq:signal} explains why. Only 1.5\% (Qwen3) and 10.9\%
(Qwen3.5) of the groups had any reward variance. If the four outcomes of a
prompt were independent draws at the mean success rate (0.80 and 0.955),
\eqref{eq:signal} would predict 59\% and 17\%. The Qwen3 student is far more
deterministic than that: when it fails a prompt, it fails all four samples.
Its entropy collapsed in stage~2 (from 0.14 to 0.002 nats per token) and
stayed below 0.05; the Qwen3.5 student stayed at 0.13--0.16 from stage~2 on. Reverse
KL predicts this: at its minimum the student's next-token distribution,
entropy included, equals the teacher's, so a confident teacher makes a
confident student (we did not measure the teachers' entropy), and a
confident student gives GRPO nothing to compare. Stage~5 of the Qwen3 run,
GRPO alone with $G=8$ and no teacher, confirmed it: after the first batch
every group had zero variance, the loss and gradient were exactly zero, and
we stopped the run at about step~20.
\end{findbox}

\paragraph{Cost.}
The Qwen3.5 run took 7.1 hours against 2.8 (about 50 Colab compute units
against 23). A token costs the same in both runs: in stage~4, a step takes
0.17 and 0.16\,s per token of mean episode length. But the Qwen3.5 student
writes episodes 2--5 times longer, with a sentence of explanation around
every tool call (stage~2: 371 against 65 tokens; stage~3: 240 against 113).

\section{Takeaways}

\begin{tcbraster}[raster columns=2, raster equal height,
    raster column skip=4mm, raster row skip=4mm]
\begin{card}{wr-terracotta}{1\enspace Distillation is the workhorse}
Prompts alone, with no labels and no reward, teach a 1.7B or 2B student to
use tools: under two hours per stage on one A100, most of the gain in the
first 50 steps.
\end{card}
\begin{card}{wr-ochre}{2\enspace RL needs tasks at the edge of ability}
GRPO learns only from prompts whose outcome varies across samples.
Distillation removes the randomness and our tasks are mostly solved; harder
tasks, larger groups or a higher temperature are the levers.
\end{card}
\begin{card}{wr-wine}{3\enspace Budgets are per model}
Same configs, but Qwen3.5 wrote episodes 2--5 times longer: limits that
never bound for Qwen3 cost it accuracy, and its run took 2.5 times the GPU
time.
\end{card}
\begin{card}{wr-sage}{4\enspace The tokenizer decides who can teach}
The teacher must share the student's vocabulary token for token, so each
family distils within itself.
\end{card}
\end{tcbraster}

\Needspace{7\baselineskip}
\paragraph{Reproduction.}
Stages 1--4 run with
\begin{codebox}
BACKEND=trl EXTRA=-{}-use\_vllm scripts/run\_stages.sh a100 1 4
\end{codebox}
\noindent The trainer is \path{opd/trl_train.py} and the configs are in
\path{opd/configs/}. Every logged metric, sampled episode and adapter is in
\path{results/trl/a100/} (Qwen3) and \path{results/trl/a100-qwen3.5/}
(Qwen3.5).

\begin{thebibliography}{99}
\setlength{\itemsep}{1pt}
\small

\bibitem[Agarwal et~al.(2024)]{agarwal2024gkd}
R.~Agarwal, N.~Vieillard, Y.~Zhou, P.~Stanczyk, S.~Ramos, M.~Geist, and
O.~Bachem.
\newblock On-policy distillation of language models: Learning from
self-generated mistakes.
\newblock In \emph{ICLR}, 2024.

\bibitem[Cobbe et~al.(2021)]{cobbe2021gsm8k}
K.~Cobbe, V.~Kosaraju, M.~Bavarian, M.~Chen, H.~Jun, L.~Kaiser, M.~Plappert,
J.~Tworek, J.~Hilton, R.~Nakano, C.~Hesse, and J.~Schulman.
\newblock Training verifiers to solve math word problems.
\newblock arXiv:2110.14168, 2021.

\bibitem[Gu et~al.(2024)]{gu2024minillm}
Y.~Gu, L.~Dong, F.~Wei, and M.~Huang.
\newblock {MiniLLM}: Knowledge distillation of large language models.
\newblock In \emph{ICLR}, 2024.

\bibitem[Hu et~al.(2022)]{hu2022lora}
E.~J. Hu, Y.~Shen, P.~Wallis, Z.~Allen-Zhu, Y.~Li, S.~Wang, L.~Wang, and
W.~Chen.
\newblock {LoRA}: Low-rank adaptation of large language models.
\newblock In \emph{ICLR}, 2022.

\bibitem[Kim and Rush(2016)]{kim2016seqkd}
Y.~Kim and A.~M. Rush.
\newblock Sequence-level knowledge distillation.
\newblock In \emph{EMNLP}, 2016.

\bibitem[Kwon et~al.(2023)]{kwon2023vllm}
W.~Kwon, Z.~Li, S.~Zhuang, Y.~Sheng, L.~Zheng, C.~H. Yu, J.~E. Gonzalez,
H.~Zhang, and I.~Stoica.
\newblock Efficient memory management for large language model serving with
{PagedAttention}.
\newblock In \emph{SOSP}, 2023.

\bibitem[Lu and Thinking Machines Lab(2025)]{lu2025onpolicy}
K.~Lu and Thinking Machines Lab.
\newblock On-policy distillation.
\newblock \emph{Thinking Machines Lab: Connectionism}, 2025.
\newblock doi:10.64434/tml.20251026.

\bibitem[Qwen Team(2026)]{qwen2026qwen35}
Qwen Team.
\newblock {Qwen3.5} model cards.
\newblock Hugging Face, 2026.

\bibitem[Ross and Bagnell(2010)]{ross2010efficient}
S.~Ross and J.~A. Bagnell.
\newblock Efficient reductions for imitation learning.
\newblock In \emph{AISTATS}, 2010.

\bibitem[Ross et~al.(2011)]{ross2011dagger}
S.~Ross, G.~J. Gordon, and J.~A. Bagnell.
\newblock A reduction of imitation learning and structured prediction to
no-regret online learning.
\newblock In \emph{AISTATS}, 2011.

\bibitem[Schulman and Thinking Machines Lab(2025)]{schulman2025lora}
J.~Schulman and Thinking Machines Lab.
\newblock {LoRA} without regret.
\newblock \emph{Thinking Machines Lab: Connectionism}, 2025.

\bibitem[Shao et~al.(2024)]{shao2024deepseekmath}
Z.~Shao, P.~Wang, Q.~Zhu, R.~Xu, J.~Song, X.~Bi, H.~Zhang, M.~Zhang,
Y.~K. Li, Y.~Wu, and D.~Guo.
\newblock {DeepSeekMath}: Pushing the limits of mathematical reasoning in open
language models.
\newblock arXiv:2402.03300, 2024.

\bibitem[von Werra et~al.(2020)]{vonwerra2020trl}
L.~von Werra, Y.~Belkada, L.~Tunstall, E.~Beeching, T.~Thrush, N.~Lambert,
S.~Huang, K.~Rasul, and Q.~Gallou{\'e}dec.
\newblock {TRL}: Transformer reinforcement learning.
\newblock \url{https://github.com/huggingface/trl}, 2020.

\bibitem[Williams(1992)]{williams1992reinforce}
R.~J. Williams.
\newblock Simple statistical gradient-following algorithms for connectionist
reinforcement learning.
\newblock \emph{Machine Learning}, 8:229--256, 1992.

\bibitem[Yang et~al.(2025a)]{yang2025qwen3}
A.~Yang et~al.
\newblock {Qwen3} technical report.
\newblock arXiv:2505.09388, 2025a.

\bibitem[Yang et~al.(2025b)]{yang2025gateddeltanet}
S.~Yang, J.~Kautz, and A.~Hatamizadeh.
\newblock Gated delta networks: Improving {Mamba2} with delta rule.
\newblock In \emph{ICLR}, 2025b.

\bibitem[Yu et~al.(2025)]{yu2025dapo}
Q.~Yu et~al.
\newblock {DAPO}: An open-source {LLM} reinforcement learning system at
scale.
\newblock arXiv:2503.14476, 2025.

\end{thebibliography}

\end{document}
